\chapter{Hiccups (Singultus)}

\section*{Definition}
Involuntary, spasmodic contractions of the diaphragm followed by sudden glottic closure producing the characteristic "hic" sound, mediated by a reflex arc involving the phrenic nerve and vagus nerve.

\section*{What Happens in Your Body}
The hiccup reflex arc involves afferent fibers (phrenic and vagus nerves, sympathetic chain from T6-T12), central coordinator in the brainstem (medulla, reticular formation), and efferent fibers (phrenic nerve to diaphragm, vagus nerve to glottis). Hiccups lasting >48h (persistent) or >1 month (intractable) require investigation for underlying pathology.

\section*{Causes}
\begin{itemize}
  \item Gastric distension (overeating, carbonated drinks, swallowing air)
  \item Sudden temperature changes (hot then cold food/drink)
  \item Alcohol consumption
  \item Excitement or emotional stress
  \item Gastroesophageal reflux disease (GERD)
  \item Medication side effects (benzodiazepines, barbiturates, corticosteroids)
  \item Central nervous system disorders (stroke, tumor, multiple sclerosis)
  \item Diaphragmatic irritation (pneumonia, pleurisy, pericarditis)
  \item Metabolic disorders (uremia, diabetes, electrolyte imbalance)
  \item Post-operative (especially abdominal or thoracic surgery)
\end{itemize}

\section*{When to See a Doctor}
See your doctor if:
\begin{itemize}
  \item Symptoms are severe or getting worse over time
  \item Hiccups persisting more than 48 hours, hiccups with severe pain, or associated with neurological symptoms
  \item You have other concerning symptoms such as fever, weight loss, or bleeding
\end{itemize}

\section*{Related Symptoms}
\begin{itemize}
  \item GERD
  \item Abdominal distension
  \item Chest discomfort
  \item Nausea
  \item Fatigue
\end{itemize}
\textbf{Important:} If this symptom persists, your doctor will check for serious underlying causes.

\section*{Self-Care / Home Management}
\begin{itemize}
  \item Try simple physical manoeuvres to interrupt the reflex arc: holding your breath, drinking a glass of cold water slowly, or swallowing a teaspoon of granulated sugar.
  \item Avoid common triggers such as carbonated beverages, spicy foods, rapid eating, and sudden temperature changes in food or drink.
  \item If hiccups persist beyond a few minutes, focus on slow, rhythmic breathing into a paper bag (not plastic) to increase CO\textsubscript{2} and reduce diaphragm spasms.
  \item Seek medical attention if hiccups last longer than 48 hours, as persistent hiccups may indicate an underlying neurological or metabolic cause.
\end{itemize}

\section*{How Your Doctor Will Evaluate You}
\begin{itemize}
  \item History focusing on onset, duration (acute vs persistent >48h vs intractable >1 month), triggers, medications, and associated neurological or gastrointestinal symptoms.
  \item Neurological examination to assess cranial nerve function, particularly phrenic and vagus nerves, and to screen for brainstem pathology.
  \item Basic metabolic panel to check for electrolyte imbalances, renal function (uraemia), and serum glucose; liver function tests if hepatic cause suspected.
  \item CT or MRI of the brain and chest is indicated for persistent or intractable hiccups to exclude structural lesions, diaphragmatic irritation, or mediastinal pathology.
\end{itemize}

\section*{Treatment Options}
\begin{itemize}
  \item Acute benign hiccups resolve spontaneously or with simple physical manoeuvres and require no specific treatment.
  \item Persistent hiccups (>48h) may be treated pharmacologically with chlorpromazine (first-line), metoclopramide, baclofen, or gabapentin under medical supervision.
  \item Treating the underlying cause (GERD with proton pump inhibitors, electrolyte repletion, discontinuing offending medications) often resolves secondary hiccups.
  \item For intractable hiccups, phrenic nerve block or vagus nerve stimulation may be considered by a specialist in refractory cases.
  \item Medications, lifestyle modifications, and physical therapies may be prescribed as appropriate.
  \item Follow-up ensures treatment effectiveness and allows timely adjustments to the care plan.
\end{itemize}

\section*{References}

\begin{thebibliography}{99}
\bibitem{harrison-hiccup} Longo DL, et al. \emph{Harrison's Principles of Internal Medicine}. 21st ed. McGraw-Hill; 2022. Chapter 50: Hiccups.
\bibitem{uptodate-hiccup} Steger M, et al. Hiccups in adults. In: UpToDate, Post TW (Ed), UpToDate, Waltham, MA; 2025.
\bibitem{straus} Straus C, et al. Hiccups: pathophysiology and management. \emph{Chest}. 2022;162(5):1108-1118.
\bibitem{kohse} Kohse EK, et al. Persistent and intractable hiccups: a systematic review. \emph{Neurology}. 2023;100(8):e825-e835.
\bibitem{wojcieszek} Wojcieszek J, et al. Chronic hiccups: etiology and treatment. \emph{Am Fam Physician}. 2023;107(2):172-179.
\bibitem{petroianu} Petroianu GA, et al. The neuroanatomy of hiccups: a review. \emph{Clin Anat}. 2022;35(4):455-464.
\end{thebibliography}
